% CLAUDE REPORT::BRIDGE-0001::2026-11-04
\documentclass{article}
\usepackage{amsmath,amssymb,amsthm}
\begin{document}
\section*{BRIDGE-0001: worst-case derivation cost, and what the bridge $SAT\in P\Rightarrow\overline W_R\le \mathrm{poly}$ can and cannot carry}
Status: PROVED/EXACT (Prop.~1, elementary); COMPUTATIONAL (tree-like table); KNOWN THEOREM (Cook--Reckhow, Haken); OPEN: superpolynomial $\overline W_R$ for a system $R$ for which the bridge is provable, and P versus NP. Nothing here changes any schedule or allocation; it adds to the shared worst-case directive (see \texttt{research/governance} Rmix341 line of LSPV-1/2/4).

\paragraph{Setup.} $R$ a Cook--Reckhow system (poly-time verifier). $W_R(F)=\min_\pi |\pi|$ over $R$-refutations (strategy search free). $\overline W_R(N)=\max_{|F|\le N}W_R(F)$. $\mathrm{LB}(R)$: $\overline W_R$ superpolynomial. $\mathrm{Br}(R)$: $SAT\in P\Rightarrow \overline W_R\le\mathrm{poly}$.

\paragraph{Proposition 1 (no-free-lunch for the bridge).} For every Cook--Reckhow $R$:
(a) $\mathrm{Br}(R)\iff(\mathrm P=\mathrm{NP}\Rightarrow R\text{ polynomially bounded})$.
(b) $\mathrm{LB}(R)\wedge\mathrm{Br}(R)\Rightarrow \mathrm P\ne \mathrm{NP}$.
(c) $\mathrm{LB}(R)\wedge \mathrm P\neq\mathrm{NP}\Rightarrow \mathrm{Br}(R)$ (vacuous truth).
Hence $\mathrm{LB}(R)\wedge\mathrm{Br}(R)\iff \mathrm{LB}(R)\wedge \mathrm P\ne\mathrm{NP}$: the pair is exactly as strong as $\mathrm P\neq\mathrm{NP}$ plus a lower bound. (d) If $\mathrm{LB}(R)$ is already a theorem (Resolution, Haken 1985), $\mathrm{Br}(R)\iff \mathrm P\ne\mathrm{NP}$; so the bridge for a fixed weak $R$ is not a lemma, it is the conclusion.
\emph{Proof.} (a) $SAT\in P\iff \mathrm P=\mathrm{NP}$ and ``polynomially bounded'' is $\overline W_R\le\mathrm{poly}$. (b) contrapose: $\mathrm P=\mathrm{NP}$ gives bounded, contradicting LB. (c) antecedent false. (d) with LB true, $\mathrm{Br}$ says $\neg(SAT\in P)$. \qed

\paragraph{Consequence for the program.} The only non-circular content is a \emph{constructive} transfer: for every correct polynomial-time decider $A$, polynomial $R$-proofs of $\mathrm{Ref}_A(F)$ (this is B1 in the LSPV-1/377 notation), obtained without assuming $\mathrm P\ne\mathrm{NP}$. This requires $R$ to formalise instancewise soundness of $A$ (propositional translation of a reflection principle). Whether truth of $\mathrm P=\mathrm{NP}$ alone gives such short proofs in Extended Frege is a bounded-arithmetic independence question (Krajicek line); I do not claim an answer. Cook--Reckhow only gives that \emph{some} polynomially bounded system exists iff $\mathrm{NP}=\mathrm{coNP}$, not that a prescribed $R$ is one. A system is a legitimate target for the bridge only if it is optimal-like; existence of $p$-optimal systems is itself open.

\paragraph{Computational table (exact min tree-like Resolution size, leaves).} Script \texttt{bridge/treelike\_min.py} (runs $<2$s, no chunking needed), data \texttt{bridge/treelike\_min\_results.json}: $PHP^{2}_{1}\!:3$, $PHP^{3}_{2}\!:11$, $PHP^{4}_{3}\!:43$; Tseitin$(K_3)\!:6$, $(K_4)\!:28$, $(K_5)\!:232$. Tree-like only (a restriction of DAG-like); finite data, no asymptotic claim; the $K_n$ growth is consistent with the known $2^{\Omega(n)}$-type tree-like Tseitin bounds but proves nothing.

\paragraph{Literature check (2026-11-04).} Search found no unrestricted Frege/EF superpolynomial lower bound. Closest: tree-like semantic Frege with bounded line size, superpolynomial (ECCC TR26-078, arXiv 2604.28172); algebraic-side EF-simulating system lower bound for the binary value principle (Theory of Computing 22, 2026). Both are restricted systems and say nothing on the bridge. Transfer classification: these are ANALOGY needing a new transfer theorem (no bridge exists from restricted tree-like semantic systems to $\mathrm{Rmix341}$).

\paragraph{Next target.} Pick the weakest $R\supseteq$ ER for which an explicit reflection-style translation of one concrete poly-time SAT-style algorithm class (e.g.\ a fixed bounded-treewidth DP, whose correctness \emph{is} provable) yields $R$-proofs of size $\mathrm{poly}$ on that class; then ask which hard families escape the class. This tests the bridge mechanism on a case where it provably holds.

\paragraph{Corpus audit (condensed).} Surveyed repository tree; read AGENTS.md, LSPV-1/377 (blob in repo), Claude 2026-10-28 notes. Duplication: Prop.~1 is elementary and implicit in B1/B2/B3 of LSPV-1/377; the explicit equivalence (c)/(d) is newly stated here, not claimed novel in the literature. No root \texttt{main.tex} ownership region belongs to this agent; \texttt{main.tex} not modified. P versus NP remains OPEN.
% END_OF_VERIFIED_REPORT::CLAUDE::BRIDGE-0001
\end{document}
